\chapter{Progetti Filtri FIR} \label{ch: Progetti Filtri FIR}

\begin{figure}[h]
    \centering
    \includegraphics[scale = 0.3]{fir digital filter.jpg}
\end{figure}  

\newpage 

\section{Progetto di filtri FIR}

\footnote{ 
    Slide del prof | 10 - Progetto Filtri FIR 1 | pag 1 - 2\\
    Appunti | 2025-11-18 | pag 1    
}

I filtri FIR, Finite Impulse Response (in Italiano, filtri a risposta impulsiva finita), rispetto agli IIR, sono filtri più semplici da progettare (proprio perchè i coefficienti del filtro sono finiti), 
ma impiegano più tempo per elaborare il segnale di ingresso $x \left[ n\right]$. \newline 

Come studiato in \ref{subsec: Sistemi causali con fase generalizzata lineare} \nameref{subsec: Sistemi causali con fase generalizzata lineare} a pag \pageref{subsec: Sistemi causali con fase generalizzata lineare}, 
se i filtri FIR sono simmetrici, allora il ritardo di gruppo è costante. \newline 

Il ritardo di gruppo, a livello di progettazione del filtro, è una caratteristica molto positiva, ma implica che il filtro avrà bisogno del tempo per andare a regime. \newline 

Come al solito, il tempo per andare a regime, chiamato anche transitorio, dipende sempre dall'applicazione in cui il filtro è applicato: 
un filtro impiegato nell'audio avrà un transitorio diverso rispetto ad un filtro impiegato nelle telecomunicazioni. \newline 

Esistono diversi metodi di progettazione di un filtro FIR, ma noi ne andremo ad analizzare 4 (di cui solo i primi 2 saranno approfonditi in questo corso, gli altri 2 saranno approfonditi alla magistrale): 

\begin{enumerate}
    \item metodo delle finestre (chiamato in inglese "Windowing"), proprio perchè è basato sulle finestre, quindi semplice, poco flessibile, non ottimo 
    \item metodo "Equiripple", grazie all'algoritmo di Parks-McClellan, è possibile progettare un filtro FIR in un modo complicato, ma flessibile e ottimo 
    \item metodo del campionamento in frequenza, il quale è un metodo mediamente complicato, non ottimo, ma flessibile 
    \item metodi di stima parametrica, metodo ottimo in generale ma ha complessità variabile 
\end{enumerate} 

\newpage 

\section{Progetto di filtri FIR: Windowing} \label{sec: Progetto di filtri FIR: Windowing}

\footnote{ 
    Slide del prof | 10 - Progetto Filtri FIR 1 | pag 3 - 4\\
    Appunti | 2025-11-18 | pag 1 \\ 
    Latex | 2025-12-16 Mattino | pag 251 
}

Prima di procedere con la progettazione di un filtro FIR, 
dobbiamo svolgere alcune osservazioni in merito ai filtri in generale. \newline 

La prima osservazione è che un filtro FIR, nel dominio z, ha la seguente funzione di rete: 

{
    \Large 
    \begin{equation}
        H \left(z\right) 
        = 
        \sum_{k = 0}^{M} 
        b_k \cdot z^{-k}
    \end{equation}
}

dove: 

\begin{itemize}
    \item i coefficienti $b_k$ del filtro FIR sono M+1 
    \item i coefficienti $b_k$ del filtro FIR sono finiti 
    \item i coefficienti $b_k$ del filtro FIR sono i campioni della risposta impulsiva, cioè di $h \left[ n\right]$ 
\end{itemize} 

La seconda osservazione è che i filtri ideali $H_d \left(e^{\jmath \omega}\right)$, cioè quelli che hanno una finestra rettangolare in frequenza, hanno risposta impulsiva $h_d \left[ n\right]$ che è IIR, cioè infinita, 
dovuto alle proprietà di Fourier tra tempo e frequenza. \newline 

\begin{tcolorbox}
    Il professor Franco Chiaraluce insegna in merito a questo argomento: le basi dell'esame di Teoria dei Segnali. \newline 

    Senza sapere le proprietà tra tempo e frequenza, l'esame di TdS è meglio non farlo. 
\end{tcolorbox} 

Grazie a questa ultima osservazione, l'idea è quella di ottenere un filtro FIR troncando la risposta impulsiva IIR dei filtri ideali calcolandoli analiticamente (in forma chiusa). \newline 

\begin{tcolorbox}
Calcolarsi i filtri in forma chiusa si intende che si calcoleranno analiticamente dei filtri, che saranno poi ottimizzati attraverso delle tecniche di ottimizzazione      
\end{tcolorbox}

Ad esempio, considerando un filtro passa-basso ideale $H_d \left(e^{\jmath \omega}\right)$ come il seguente: 

{
    \Large 
    \begin{equation}
        H_d \left(e^{\jmath \omega}\right) 
        = 
        \begin{cases}
            1 \quad &\text{per }\abs{\omega} < \omega_c 
            \\ 
            0 \quad &\text{altrove } 
        \end{cases}
    \end{equation}
}

dove $\omega_c$ è la frequenza di cut-off del filtro desiderato. \newline 

Graficato nello spettro, $H_d \left(e^{\jmath \omega}\right)$ avrà il seguente andamento: 

\begin{figure}[h]
    \centering
    \includegraphics[scale = 0.15]{finestra rettangolare in frequenza.png}
\end{figure}  

\newpage 

Anti-trasformando in Fourier $H_d \left(e^{\jmath \omega}\right)$, cioè la finestra rettangolare, avremo $h_d \left[ n\right]$, 
che vale: 

{
    \Large 
    \begin{equation}
        \begin{split}
        H_d \left(e^{\jmath \omega}\right) 
        &= 
        \begin{cases}
            1 \quad &\text{per }\abs{\omega} < \omega_c 
            \\ 
            0 \quad &\text{altrove } 
        \end{cases} 
        \\ 
        &\downarrow 
        \\ 
        h_d \left[ n\right] 
        &= 
        \dfrac{\sin \left(\omega_c n\right)}{\pi n}
        \end{split}
    \end{equation}
}

In altre parole, da finestra rettangolare in Fourier, $h_d \left[ n\right]$, nel dominio n, è un sinc, come mostrato in figura: 

\begin{figure}[h]
    \centering
    \includegraphics[scale = 0.8]{Sinc nel dominio n.png}
\end{figure} 

e che quindi $h_d \left[ n\right]$ si estende per $n \in \left[ - \infty, + \infty\right]$, quindi il sinc è un filtro IIR. \newline 

Allora, per non avere un filtro che si estende tra $- \infty$ e $+ \infty$, possiamo troncare $h_d \left[ n\right]$, 
come si vede nella seguente figura: 

\begin{figure}[h]
    \centering
    \includegraphics[scale = 0.8]{Sinc troncato non causale.png}
\end{figure} 

Bene aver troncato $h_d \left[ n\right]$, ma rimane un filtro non causale, 
perchè il filtro $h_d \left[ n\right]$ non parte dall'indice $n = 0$. \newline 

\newpage 

Per renderlo causale, basta shiftarlo (cioè ritardarlo), come si vede in figura: 

\begin{figure}[h]
    \centering
    \includegraphics[scale = 0.8]{sinc troncato e causalizzato.png}
\end{figure} 

Ora $h_d \left[ n\right]$ è un sinc, troncato e causalizzato, quindi è un filtro FIR pronto da essere utilizzato dal calcolatore. \newline 

Quindi in formule siamo partiti da $H_d \left(e^{\jmath \omega}\right)$: 

{
    \Large 
    \begin{equation}
        H_d \left(e^{\jmath \omega}\right) 
        = 
        \begin{cases}
            1 \quad &\text{per } \abs{\omega} < \omega_c 
            \\ 
            0 \quad &\text{altrove } 
        \end{cases} 
    \end{equation}
} 

Abbiamo fatto la sua anti-trasformata di Fourier: 

{
    \Large 
    \begin{equation}
        h_d \left[ n\right] 
        = 
        \dfrac{1}{2 \pi}
        \int_{- \pi}^{+ \pi} 
        H_d \left(e^{\jmath \omega}\right) 
        \cdot 
        e^{\jmath \omega n} 
        d \omega
    \end{equation}
} 

Una volta shiftata $h_d \left[ n\right]$ troncata per n-indici M+1, l'abbiamo traslata e fatta partire dall'indice $n = 0$, 
per rendere $h_d \left[ n\right]$ causale. \newline 

Alla fine di tutto ciò, possiamo definire la risposta impulsiva $h \left[ n\right]$ del filtro sinc, troncato e causalizzato FIR come: 

{
    \Large 
    \begin{equation}
        h \left[ n\right] 
        = 
        \begin{cases}
            h_d \left[ n\right] \quad &\text{per } 0 \leq n \leq M
            \\ 
            0 \quad &\text{altrove } 
        \end{cases}
    \end{equation}
}

Ma, definendo una nuova finestra rettangolare $w \left[ n\right]$ nel dominio n: 

{
    \Large 
    \begin{equation}
        w \left[ n\right] 
        = 
        \begin{cases}
            1 \quad &\text{per } 0 \leq n \leq M
            \\ 
            0 \quad &\text{altrove } 
        \end{cases}
    \end{equation}
}

possiamo definire la risposta impulsiva del filtro $h \left[ n\right]$ anche come: 

{
    \Large 
    \begin{equation}
        h \left[ n\right] 
        = 
        h_d \left[ n\right] 
        \cdot 
        w \left[ n\right]
    \end{equation}
}

\newpage 

\subsection{Conseguenze del Windowing} \label{subsec: Conseguenze del Windowing}

\footnote{ 
    Slide del prof | 10 - Progetto Filtri FIR 1 | pag 5 - 6\\
    Appunti | 2025-11-18 | pag 1 - 2 
} 

Analizzando $H \left(e^{\jmath \omega}\right)$, ma questa volta con il windowing $w \left[ n\right]$: 

{
    \Large 
    \begin{equation}
        \setlength{\jot}{10pt} % affecting the line spacing in the environment
        \begin{split}
            H \left(e^{\jmath \omega}\right) 
            &= 
            DTFT \left[ h \left[ n\right]\right] 
            \\
            &= 
            DTFT \left[ h_d \left[ n\right] \cdot w \left[ n\right]\right] 
            \\
            &= 
            \dfrac{1}{2 \pi} 
            \int_{- \pi}^{+ \pi} 
            H_d \left(e^{\jmath \theta}\right) 
            \cdot 
            W \left(e^{\jmath \left( \omega - \theta\right)}\right) 
            d \theta
        \end{split}
    \end{equation}
}

cioè, come abbiamo studiato nei corsi precedenti, un prodotto nel tempo, cioè $h_d \left[ n\right] \cdot w \left[ n\right]$, 
genere una convoluzione complessa in $\omega$, cioè $\int_{- \pi}^{+ \pi} H_d \left(e^{\jmath \theta}\right) \cdot W \left(e^{\jmath \left( \omega - \theta\right)}\right) d \theta$. \newline 

Poichè $H_d \left(e^{\jmath \omega}\right)$, cioè la finestra rettangolare in Fourier, è periodica, 
allora possiamo considerare per $h_d \left[ n\right]$ i suoi coefficienti della serie di Fourier associata (e non i coefficienti della trasformata di Fourier, cioè DTFT). \newline 

\begin{tcolorbox}
    Ricordiamo al volo che la serie di Fourier viene utilizzata per segnali che sono periodici nel tempo e si estendono per $n \left[ - \infty, + \infty\right]$. \newline 

    Invece le trasformate, quindi la DTFT, viene utilizzata per segnali finiti non periodici e che non si estendono per $n \left[ - \infty, + \infty\right]$. 
\end{tcolorbox}  

Quindi, svolgere il windowing (o anche detto finestrare) significa troncare, cioè approssimare, una serie di Fourier, come si vede in figura: 

\begin{figure}[h]
    \centering
    \includegraphics[scale = 0.2]{windowing differenza.png}
\end{figure} 

Come si vede nella figura, $H_d$ (che è in rosso) dimostra l'andamento di un filtro rettangolare ideale, 
invece H (che è in blu) dimostra l'andamento di un filtro "rettangolare" dovuto al finestramento. \newline 

Le oscillazioni, sia in banda passante che in banda oscura, sono dovute al fenomeno di Gibbs (che avevamo incontrato in \ref{subsec: Filtri passa basso ideale} \nameref{subsec: Filtri passa basso ideale} a pag \pageref{subsec: Filtri passa basso ideale}). \newline 

Inoltre, le oscillazioni aumentano all'aumentare di M e sono maggiori nella banda di transizione (cioè nella discontinuita per $H_d$). \newline 

Ricordiamo al volo che M+1 è la lunghezza della finestra $w \left[ n\right]$. \newline 

Ma la vera domanda da porsi per la progettazione dei filtri FIR è: è necessario avvicinarsi ad un filtro ideale? \newline 

La risposta è no, perchè il filtro impiegherebbe troppo tempo. \newline 

\newpage 

Inoltre, in una finestratura, abbiamo due parametri liberi per svolgere questa tecnica, che sono: 

\begin{itemize}
    \item la lunghezza della finestra, cioè quanti M+1 campioni servono al filtro FIR 
    \item la forma della finestra (non per forza deve essere rettangolare come abbiamo fatto negli esempi precedenti)
\end{itemize} 

Considerando, per adesso e per semplicità, una finestra $w \left[ n\right]$ nel tempo rettangolare, in Fourier sarà uguale a: 

{
    \Large 
    \begin{equation}
        \setlength{\jot}{10pt} % affecting the line spacing in the environment
        \begin{split}
            W \left(e^{\jmath \omega}\right) 
            &= 
            \sum_{n = 0}^{M} 
            w \left[ n\right] 
            \cdot 
            e^{- \jmath \omega n} 
            \\
            &= 
            \sum_{n = 0}^{M} 
            1 
            \cdot 
            e^{- \jmath \omega n} 
            \\
            &= 
            \sum_{n = 0}^{M} 
            e^{- \jmath \omega n} 
            \\
            &= 
            \dfrac{1 - e^{- \jmath \omega \left(M + 1\right)}}{1 - e^{- \jmath \omega}} 
            \\
            &= 
            e^{- \jmath \omega \frac{M}{2}} 
            \cdot 
            \dfrac{\sin \left[ \dfrac{\omega \left(M + 1\right)}{2}\right]}{\sin \left[ \dfrac{\omega}{2}\right]}
        \end{split}
    \end{equation}
}

Cioè, la finestra rettangolare nel tempo $w \left[ n\right]$, diventa, in Fourier, un rapporto tra due seni, come mostrato in figura: 

\begin{figure}[h]
    \centering
    \includegraphics[scale = 0.35]{Finestra rettangolare nel dominio n nello spettro.png}
\end{figure} 

Proprio perchè nella formula di $W \left(e^{\jmath \omega}\right) $ è presente il fatto M, 
possiamo svolgere queste osservazioni, sapendo che l'area dei lobi è costante, all'aumentare di M: 

\begin{itemize}
    \item il lobo principale si stringe 
    \item aumentano i lobi laterali 
\end{itemize} 

\newpage 

\subsection{Esigenze lunghezza finestra}

\footnote{ 
    Slide del prof | 10 - Progetto Filtri FIR 1 | pag 7\\
    Appunti | 2025-11-18 | pag 2  
} 

La scelta della lunghezza della finestra, cioè di quanti M+1 indici deve essere lunga $w \left[ n\right]$ nel dominio n (la cui trasformata di Fourier è $W \left(e^{\jmath \omega}\right)$), 
è cruciale perchè al variare di M si possono fare queste seguenti osservazioni: 

\begin{itemize}
    \item $w \left[ n\right]$ corta implica minore complessità del filtro 
    \item $W \left(e^{\jmath \omega}\right)$ stretta implica minore distorsione nel filtro finale $H \left(e^{\jmath \omega}\right)$
\end{itemize}

\begin{tcolorbox}
    La scelta della lunghezza della finestra è un altro problema ingegneristico che prende il nome di "la coperta è corta" (come al solito), o come dicono gli economisti "non c'è budget". 
\end{tcolorbox}

Come abbiamo scritto precedentemente, il filtro finale $H \left(e^{\jmath \omega}\right)$ sarà la convoluzione tra il filtro ideale rettangolare $H_d \left(e^{\jmath \omega}\right)$ e la trasformata della finestra $W \left(e^{\jmath \omega}\right)$: 

\begin{figure}[h]
    \centering
    \includegraphics[scale = 0.35]{windowing convoluzione esempio 1.png}
\end{figure} 

Invece, aumentando M possiamo avere il seguente caso: 

\begin{figure}[h]
    \centering
    \includegraphics[scale = 0.35]{Windowing con M maggiore.png}
\end{figure} 

possiamo notare che: 

\begin{itemize}
    \item la banda di transizione dipende dalla larghezza del main lobe (in italiano lobo principale), quindi dal fattore M 
    \item il ripple dipende dall'ampiezza dei lobi laterali, in inglese sidelobes (ma generalmente i sidelobes non sono importanti perchè sono in banda oscura, l'importante che abbiano un'ampiezza minore delle specifiche)
\end{itemize} 

\newpage 

\subsection{Elenco delle finestre principali} \label{subsec: Elenco delle finestre principali}

\footnote{ 
    Slide del prof | 10 - Progetto Filtri FIR 1 | pag 8 - 12\\
    Appunti | 2025-11-18 | pag 2 \\ 
    Latex | 2025-12-16 Mattino | pag 259 
} 

Siccome si vuole scegliere una finestra $w \left[ n\right]$ tale che in Fourier $W \left(e^{\jmath \omega}\right)$ abbia dei lobi laterali più bassi, 
ciò è possibile azzerando dolcemente gli estremi della finestra. \newline 

A questo punto, siccome l'area sotto alla finestra è costante in tutti i lobi, il lobo principale è più largo. \newline 

Grazie a questo principio, negli anni, diversi studiosi hanno proposto diverse finestre, 
che sono nulla fuori dalla lunghezza della finestra stessa. \newline 

Di seguito le finestre $w \left[ n\right]$ più utilizzati nell'ambito DSP. \newline 

\textbf{Rettangolare} \newline 

{
    \Large 
    \begin{equation}
        w \left[ n\right] = 0 \quad \text{ per } 0 \leq n \leq M
    \end{equation}
}

\textbf{Bartlett} \newline 

{
    \Large 
    \begin{equation}
        w \left[ n\right] 
        = 
        \begin{cases}
            \dfrac{2 n}{M} \quad &\text{per } 0 \leq n \leq \dfrac{M}{2} 
            \\ 
            2 - \dfrac{2 n}{M} \quad &\text{per }\dfrac{M}{2} \leq n \leq M 
        \end{cases}
    \end{equation}
}

\textbf{Hanning} \newline 

{
    \Large 
    \begin{equation}
        w \left[ n\right] 
        = 
        0.5 - 0.5 \cos \left(\dfrac{2 \pi n}{M}\right) 
        \quad \text{ per } 0 \leq n \leq M
    \end{equation}
}

\textbf{Hamming} \newline 

{
    \Large 
    \begin{equation}
        w \left[ n\right] 
        = 
        0.54 - 0.46 \cos \left(\dfrac{2 \pi n}{M}\right) 
        \quad \text{ per } 0 \leq n \leq M
    \end{equation}
}

\begin{tcolorbox}
    Sembrano simili, ma Hammming è con la m, questa invece è Hanning è con le n: sono due finestre diverse.  
\end{tcolorbox} 

\textbf{Blackman} \newline 

{
    \Large 
    \begin{equation}
        w \left[ n\right] 
        = 
        0.42 - 0.5 \cos \left(\dfrac{2 \pi n}{M}\right) + 0.08 \cos \left(\dfrac{4 \pi n}{M}\right) 
        \quad \text{ per } 0 \leq n \leq M
    \end{equation}
}

\newpage 

Da un punto di vista grafico, le varie finistre $w \left[ n\right]$ appena elencate hanno questo andamento (considerando le curve come continue e non discrete come dovrebbero essere realmente): 

\begin{figure}[h]
    \centering
    \includegraphics[scale = 0.4]{finestre rettangoli come curve continue.png}
\end{figure} 

Come possiamo visualizzare dalla figura, tutte le finestre $w \left[ n\right]$ sono simmetriche rispetto all'indice $n = \dfrac{M}{2}$ (questa è una caratteristica molto importante nella progettazione dei filtri FIR). \newline 

Ora andiamo a confrontarle in Fourier. \newline 

Di seguito lo spettro di una finestra rettangolare: 

\begin{figure}[h]
    \centering
    \includegraphics[scale = 0.4]{finestre rettangolare nello spettro.png}
\end{figure} 

dove la larghezza del lobo principale $ \Delta \omega $ vale: 

{
    \Large 
    \begin{equation}
        \Delta \omega 
        = 
        \dfrac{4 \pi}{M + 1}
    \end{equation}
}

\newpage 

L'attenuazione minima di una finestra rettangolare generalmente si misura in decibel, cioè in dB, ed è di almeno 21 dB. \newline 

 Nella figura è indicato -13 dB, ma in realtà, in una finestra rettangolare, si considera il suo magnitude-squared: ecco perchè -21 dB. \newline 
 
Invece, una finestra di Bartlet triangolare ha il seguente spettro: 

\begin{figure}[h]
    \centering
    \includegraphics[scale = 0.4]{finestre bartlett triangolare nello spettro.png}
\end{figure} 

dove la larghezza del lobo principale $ \Delta \omega $ vale: 

{
    \Large 
    \begin{equation}
        \Delta \omega 
        = 
        \dfrac{8 \pi}{M + 1}
    \end{equation}
}

L'attenuazione minima di una finestra Bartlett triangolare è di almeno 25 dB. \newline 

\begin{tcolorbox}
Questa è l'unica finestra in cui anche nella figura è indicata l'attenuazione in magnitude-squarred.     
\end{tcolorbox}

Confrontando una finestra rettangolare con una finestra di Bartlett rettangolare, avremo un'attenuazione maggiore minima all'estremo del lobo principale 
(-25 dB è un'attuazione maggiore rispetto ai -21 dB), ma a discapito di un allargamento del lobo principale $\Delta \omega $ (si passa da $\dfrac{4 \pi}{M + 1}$ a $\dfrac{8 \pi}{M + 1}$ ). \newline 

\newpage 

Invece, una finestra di Hanning ha il seguente spettro: 

\begin{figure}[h]
    \centering
    \includegraphics[scale = 0.4]{finestre hanning nello spettro.png}
\end{figure} 

dove la larghezza del lobo principale $ \Delta \omega $ vale: 

{
    \Large 
    \begin{equation}
        \Delta \omega 
        = 
        \dfrac{8 \pi}{M}
    \end{equation}
}

L'attenuazione minima di una finestra di Hanning è di almeno 44 dB. \newline  

Invece, una finestra di Hamming ha il seguente spettro: 

\begin{figure}[h] \label{fig: finestra di Hamming}
    \centering
    \includegraphics[scale = 0.4]{finestre hamming nello spettro.png}
\end{figure} 

dove la larghezza del lobo principale $ \Delta \omega $ vale: 

{
    \Large 
    \begin{equation}
        \Delta \omega 
        = 
        \dfrac{8 \pi}{M}
    \end{equation}
}

L'attenuazione minima di una finestra di Hanning è di almeno 53 dB. \newline  

Confrontando una finestra di Hanning con una finestra di Hamming, avremo un'attenuazione maggiore minima all'estremo del lobo principale per una finestra di Hanning 
(-53 dB è un'attuazione maggiore rispetto ai -44 dB), con la stessa larghezza del lobo principale $\Delta \omega $ (che è di $\dfrac{8 \pi}{M}$ ). \newline  

L'ultima, ma non per importanza, abbiamo la finestra Blackman, che ha il seguente andamento nello spettro: 

\begin{figure}[h]
    \centering
    \includegraphics[scale = 0.4]{finestre blackman nello spettro.png}
\end{figure} 

dove la larghezza del lobo principale $ \Delta \omega $ vale: 

{
    \Large 
    \begin{equation}
        \Delta \omega 
        = 
        \dfrac{12 \pi}{M}
    \end{equation}
}

L'attenuazione minima di una finestra di Hanning è di almeno 74 dB. \newline  

Rispetto alle finestre elencate, la Hamming è la migliore finestra per quanto riguarda l'attenuazione, proprio perchè attenua di 74 dB, 
ma è quella che ha una larghezza maggiore del lobo principale. \newline 

\newpage 

\subsection{Caratteristiche del metodo Windowing}

\footnote{ 
    Slide del prof | 10 - Progetto Filtri FIR 1 | pag 13 \\
    Appunti | 2025-11-18 | pag 2  
} 

\begin{tcolorbox}
Queste sono le considerazioni che realmente sono necessarie da memorizzare per l'esame.     
\end{tcolorbox}

Dopo tutte queste spiegazioni, come mostrato dalla seguente figura:     

\begin{figure}[h]
    \centering
    \includegraphics[scale = 0.35]{metodo windowing scelta della finestra.png}
\end{figure} 

dove: 

\begin{itemize}
    \item $H_d \left(e^{\jmath \omega}\right)$ è il filtro ideale 
    \item $H \left(e^{\jmath \omega}\right)$ è il filtro realizzato con tecnica windowing 
    \item $W \left(e^{\jmath \omega}\right)$ è la finestra troncata 
\end{itemize} 

Quando si utilizza la tecnica Windowing per la progettazione di un filtro FIR bisogna ricordarsi delle seguente caratteristiche di questa tecnica: 

\begin{itemize}
    \item la largezza della banda di transione è inversamente proporzionale ad M e dipende dal tipo di finestra (cioè se si usa una finestra rettangolare, è una di Bartlett, o una Hanning, ecc..) 
    \item il ripple, cioè l'attenuazione minima in banda oscura, dipende solo dal tipo di finestra, e non dall'M scelto  
\end{itemize} 

Riassumendo, nella tecnica Windowing, una volta scelta la finestra, si calcola $\Delta \omega$, cioè la larghezza del lobo principale, da cui poi si può ricavare il coefficiente M. \newline 

\newpage 

\subsection{Finestre di Kaiser}

\footnote{ 
    Slide del prof | 10 - Progetto Filtri FIR 1 | pag 14 - 15 \\
    Appunti | 2025-11-18 | pag 2 - 3
}  

Kaiser nel 1974 ha proposto una famiglia flessibile di finestre di tipo sub-ottimo, cioè in cui la maggior parte dell'energia si trova nel lobo principale per una data ampiezza di picco 
del lobo laterale più vicino. \newline 

La finestra di Kaiser ha la seguente formula: 

{
    \Large 
    \begin{equation}
        w \left[ n\right] 
        = 
        \dfrac{
            I_0 \left[ 
                \beta 
                \cdot 
                \sqrt{
                    1 
                    - 
                    \left[ 
                        \dfrac{n - \alpha}{\alpha}
                    \right]^{2}
                }
            \right]
        }{I_0 \left[ \beta\right]}
    \end{equation}
}


dove: 

\begin{itemize}
    \item $\alpha$ è un parametro e vale $\alpha = \dfrac{M}{2}$ 
    \item la finestra $w \left[ n\right]$ si estende per gli indici $0 \leq n \leq M$ 
    \item $\beta$ è un parametro 
    \item $I_0 \left[ \cdot\right]$ è la funzione di Bessel modificata di ordine zero del primo tipo 
\end{itemize}


\begin{tcolorbox}
    La formula di Kaiser va ricordata per l'esame. \newline 

    La funzione di Bessel l'avevamo già studiata nel corso di Campi Elettromagnetici, ma per un breve ripasso al volo: 

    \begin{itemize}
        \item Project IV: Bessel Functions and their Zeros | Lecture 47 | Numerical Methods for Engineers by Jeffrey Chasnov\\\url{https://www.youtube.com/watch?v=NoblEFn49sw} 
        \item Modified Bessel Functions by ICMEStudio\\\url{https://www.youtube.com/watch?v=_X03zgkOX0c}
    \end{itemize}

\end{tcolorbox}

La formula di $w \left[ n\right]$ ha due gradi di libertà che sono $\beta$ e M. \newline 

Come si vede dalla seguente figura: 

\begin{figure}[h]
    \centering
    \includegraphics[scale = 0.35]{finestra di Kaiser.png}
\end{figure} 

\newpage 

Se il valore di $\beta$ aumenta, avremo che: 

\begin{itemize}
    \item si allarga il lobo principale 
    \item si abbassano i lobi laterali (e se il valore di $\beta$ è nullo, allora $w \left[ n\right]$ diventa una finestra rettangolare ideale)
\end{itemize} 

Invece, all'aumentare di M: 

\begin{itemize}
    \item si stringe il lobo principale 
    \item l'altezza dei lobi laterali non cambia 
\end{itemize} 

In genere, con le finestre di Kaiser, si sceglie prima il valore di $\beta$ piuttosto che il valore di M 
perchè il valore di $\beta$ sceglie il tipo di finestra da utilizzare. \newline 

\newpage 

\subsubsection{Metodo di Kaiser} \label{subsubsec: Metodo di Kaiser}

\footnote{ 
    Slide del prof | 10 - Progetto Filtri FIR 1 | pag 16\\
    Appunti | 2025-11-18 | pag 2
}  

Il metodo di Kaiser è un metodo empirico, quindi senza dimostrazione matematica ma trovata per tentativi, 
il quale permette di calcolare i coefficienti di M e $\beta$, e quindi progettare il filtro attraverso le specifiche $\omega_p$, $\omega_s$ e $\delta$ senza il calcolo della larghezza del lobo principale $\Delta \omega$. \newline 

Il metodo di Kaiser comprende le due seguenti formule empiriche: 

{
    \Large 
    \begin{equation}
        \beta 
        = 
        \begin{cases}
        0.1102 \cdot \left(A - 8.7\right) 
        \quad &\text{per } A > 50 
        \\
        0.5842 \cdot \left(A - 21\right)^{0.4} + 0.07886 \cdot \left( A - 21\right)    
        \quad &\text{per } 21 \leq A \leq 50 
        \\
        0  
        \quad &\text{per } A < 21
        \end{cases}
    \end{equation}
} 

{
    \Large 
    \begin{equation}
        M 
        = 
        \dfrac{A - 8}{2.285 \cdot \left(\omega_s - \omega_p\right)}
    \end{equation}
} 

dove: 

\begin{itemize}
    \item il valore di M è calcolato con una incertezza di $\pm 2$ 
    \item A è il ripple al quadrato in dB e vale: 
    {
        \Large 
        \begin{equation}
            A = - 20 \cdot \log_{10} \left(\delta\right)
        \end{equation}
    }
\end{itemize}

\begin{tcolorbox}
    Ai fini dell'esame, non dobbiamo ricordare le formule di $\beta$ e M, ma dobbiamo solo ricordarci che $\beta$ varia per 3 range di A e M lega le specifiche senza calcolarsi il lobo principale. 
\end{tcolorbox}

\newpage 

\subsection{Conclusioni sul metodo windowing}

\footnote{ 
    Slide del prof | 10 - Progetto Filtri FIR 1 | pag 17 \\
    Appunti | 2025-11-18 | pag 3 
}  

Dopo tutta questa spiegazione dettagliata sul metodo Windowing per la progettazione dei filtri FIR, 
possiamo trarre le seguenti considerazioni. \newline 

Dato una qualsiasi finestra, possiamo dire che: 

\begin{itemize}
    \item il ripple $\delta$ nella banda passante e nella banda oscura è lo stesso e dipende dal tipo di finestra 
    \item la larghezza della banda di transizione è la stessa per ogni transizione e non dipende dal tipo di finestra e dalla lunghezza del filtro 
\end{itemize} 

\newpage 

\subsection{Tipi di finestre "a catalogo"}

\footnote{ 
    Slide del prof | 10 - Progetto Filtri FIR 2 | pag 2 - 3 \\
    Appunti | 2025-11-18 | pag 3 
}  

Se volessimo progettare un filtro FIR di tipo I o di tipo II potremmo utilizzare questi tipi di finestre "a catalogo" con i seguenti valori: 

\begin{figure}[h]
    \centering
    \includegraphics[scale = 0.35]{tipi di finestre caratteristiche.png}
\end{figure} 

Consideriamo solo filtri FIR di tipo I o di tipo II perchè i filtri di tipo III e di tipo IV, come studiato in 
\ref{subsec: Sistemi causali con fase generalizzata lineare} \nameref{subsec: Sistemi causali con fase generalizzata lineare} a pag \pageref{subsec: Sistemi causali con fase generalizzata lineare}, 
questi due ultimi non sono simmetrici, quindi perdono la simmetria del sinc. \newline 

Confrontando la precedente tabella con la seguente figura (con note): 

\begin{figure}[h]
    \centering
    \includegraphics[scale = 0.4]{errore delle finestre.png}
\end{figure} 

notiamo che le finestre Kaiser sono più raffinate rispetto alle altre finestre, cioè si avvicinano di più alla finestra ideale rispetto alle altre finestre. \newline 

\newpage 

\subsubsection{Esempio di progettazione di filtri FIR} \label{subsubsec: Esempio di progettazione di filtri FIR}

\footnote{ 
    Slide del prof | 10 - Progetto Filtri FIR 2 | pag 4 - 5 \\
    Appunti | 2025-11-18 | pag 3 \\ 
    Latex | 2025-12-16 Mattino | pag 270 
}  

Consideriamo un esempio di progettazione di un filtro con i calcoli fatti "carta e penna", senza filterDesigner di Matlab o altri programmi. \newline 

Di seguito la richiesta. \newline 

Progettare il filtro passa-banda in figura: 

\begin{figure}[h]
    \centering
    \includegraphics[scale = 0.2]{Filtro passa-banda FIR esempio.png}
\end{figure} 

con un'attenuazione minima in banda oscura di 50 dB e bande di transizione inferiori a 0.05 (f). \newline 

Sapendo che la frequenza $f_{c1}$ e $f_{c2}$ valgono, rispettivamente: 

{
    \Large 
    \begin{equation}
        f_{c1} = 0.2 \text{ Hz} 
    \end{equation}
}

{
    \Large 
    \begin{equation}
        f_{c2} = 0.3 \text{ Hz} 
    \end{equation}
}

Allora le pulsazioni di taglio sono, rispettivamente: 

{
    \Large 
    \begin{equation}
        \begin{split}
            \omega_{c1} &= 2 \pi f_{c1} 
            \\
            &= 2 \pi \cdot 0.2 
            \\
            &= 1.25 
        \end{split} 
    \end{equation}
}

{
    \Large 
    \begin{equation}
        \begin{split}
            \omega_{c2} &= 2 \pi f_{c2} 
            \\
            &= 2 \pi \cdot 0.3 
            \\
            &= 1.88
        \end{split} 
    \end{equation}
}

La larghezza delle bande di transizioni, che sono uguali sia a destra che a sinistra, valgono: 

{
    \Large 
    \begin{equation}
        \begin{split}
        \Delta \omega 
        &= 
        \omega_{s1} - \omega_{p1} 
        = 
        \omega_{s2} - \omega_{p2} 
        \\
        &= 
        2 \pi \cdot 0.05 
        \\
        &= 
        0.31 
        \end{split}
    \end{equation}
}

\newpage 

Siccome il progetto ci permette di scegliere un'attenuazione minima, significa che possiamo scegliere qualsiasi finestra che noi vogliamo, 
l'importante che rispecchino le specifiche. \newline 

Considerando una finestra di Hamming, come scritto in \ref{fig: finestra di Hamming} \nameref{fig: finestra di Hamming} a pag \pageref{fig: finestra di Hamming}: 

{
    \Large 
    \begin{equation}
        \Delta \omega = \dfrac{8 \pi}{M}
    \end{equation}
}

Scrivendo la formula inversa, possiamo scrivere che il valore di M vale: 

{
    \Large 
    \begin{equation}
        \setlength{\jot}{10pt} % affecting the line spacing in the environment
        \begin{split}
            M 
            &= \dfrac{8 \pi}{\Delta \omega} 
            \\
            &= \dfrac{8 \pi}{0.31} 
            \\
            &\approx 82  
        \end{split} 
    \end{equation}
}

La finestra di Hamming con $M = 82$ rispetta le specifiche del progetto perchè ha un'attenuazione minima di 53 dB. \newline 

Invece, considerando una finestra Kaiser, come scritto in \ref{subsubsec: Metodo di Kaiser} \nameref{subsubsec: Metodo di Kaiser} a pag \pageref{subsubsec: Metodo di Kaiser}, 
data un'attenuazione da specifica di A uguale a: 

{
    \Large 
    \begin{equation}
        A = 50 \text{ dB}
    \end{equation}
}

abbiamo che il valore di $\beta$ (essendo A nell'intervallo $21 \leq A \leq 50$) è di: 

{
    \Large 
    \begin{equation}
        \begin{split}
            \beta 
            &= 
            0.5842 \cdot \left(A - 21\right)^{0.4} + 0.07886 \cdot \left( A - 21\right) 
            \\
            &= 
            0.5842 \cdot \left(50 - 21\right)^{0.4} + 0.07886 \cdot \left( 50 - 21\right) 
            \\
            &= 
            0.5842 \cdot \left(29\right)^{0.4} + 0.07886 \cdot \left( 29\right) 
            \\
            &= 
            4.53
        \end{split}
    \end{equation}
}

Sempre dallo stesso metodo di Kaiser, possiamo calcolarci il valore di M, che è di: 

{
    \Large 
    \begin{equation} 
        \setlength{\jot}{10pt} % affecting the line spacing in the environment
        \begin{split}
        M 
        &= 
        \dfrac{A - 8}{2.285 \cdot \left(\omega_s - \omega_p\right)} 
        \\
        &= 
        \dfrac{A - 8}{2.285 \cdot \Delta \omega} 
        \\ 
        &= 
        \dfrac{50 - 8}{2.285 \cdot 0.31} 
        \\
        &= 
        60  
        \end{split}
    \end{equation}
} 

\newpage 

Dalle formule precedenti, considerando le formule date come filtri passa-basso, si può ricavare i valori di M e della finestra $w \left[ n\right]$ per un filtro passa banda 
con risposta impulsiva $h_d \left[ n\right]$. \newline 

Si noti che, in Fourier, avremo che la risposta del filtro passa-banda $H_d \left(e^{\jmath \omega}\right)$ è di: 

{
    \Large 
    \begin{equation}
        H_d \left(e^{\jmath \omega}\right) 
        = 
        H_{LP2} \left(e^{\jmath \omega}\right) 
        - 
        H_{LP1} \left(e^{\jmath \omega}\right) 
    \end{equation}
}

in cui: 

\begin{itemize}
    \item $H_{LP2} \left(e^{\jmath \omega}\right)$ è un filtro passa-basso con pulsazione di taglio a $\omega_{c2}$ 
    \item $H_{LP1} \left(e^{\jmath \omega}\right)$ è un filtro passa-basso con pulsazione di taglio a $\omega_{c1}$
\end{itemize}

Anti-trasformando $H_d \left(e^{\jmath \omega}\right)$, avremo che la risposta impulsiva $h_d \left[ n\right]$ sarà uguale a: 

{
    \Large 
    \begin{equation}
        h_d \left[ n\right]
        = 
        2 \dfrac{\sin \left(\omega_{c2} n\right)}{\pi n} 
        - 
        2 \dfrac{\sin \left(\omega_{c1} n\right)}{\pi n} 
    \end{equation}
}

perchè l'anti-trasformata di una finestra rettangolare è proprio un sinc. \newline 

Se consideriamo un filtro a "fase zero", quindi non causale, avremo che la sua risposta impulsiva $h \left[ n\right]$ sarà uguale a: 

{
    \Large 
    \begin{equation}
        h \left[ n \right] 
        = 
        h_d \left[ n\right]
        \cdot 
        w \left[ n + \dfrac{M}{2}\right]
        \text{ per } 
        n = - \dfrac{M}{2}, \dots, 0, \dots, + \dfrac{M}{2}  
    \end{equation}
} 

Invece, se consideriamo un filtro a fase lineare, quindi causale ma con ritardo di gruppo pari a $\dfrac{M}{2}$, avremo che la sua risposta impulsiva $h \left[ n\right]$ sarà uguale a: 

{
    \Large 
    \begin{equation}
        h \left[ n\right] 
        = 
        h_d \left[ n - \dfrac{M}{2}\right] 
        \cdot 
        w \left[ n\right]
        \text{ per } 
        n = 0, \dots, M 
    \end{equation}
}

dove per $w \left[ n\right]$ è la finestra scelta. \newline 

Grazie alle due espressioni di $h \left[ n\right]$, abbiamo dimostrato che, con la tecnica windowing, si possono progettare dei filtri che sono a fase lineare perchè si mantiene la simmetria. \newline  

Inoltre, grazie alla tecnica windowing, se si finestra in modo causale, significa anche ritardare, avendo così un ritardo di gruppo costante. \newline 

È possibile svolgere lo stesso procedimento anche per un filtro di tipo II. \newline 

\newpage 

\section{Progetto di filtri FIR: Equiripple}

\footnote{ 
    Slide del prof | 10 - Progetto Filtri FIR 2 | pag 7 - 9\\
    Appunti | 2025-11-19 | pag 1 \\ 
    Latex | 2025-12-16 Mattino | pag 272 
}

Come studiato nel capitolo precedente relativo alla progettazione dei filtri IIR, 
con la tecnica windowing, si vuole progettare un filtro FIR in cui, come gli IIR: 

\begin{itemize}
    \item l'errore di approssimazione è concentrato nei dintorni della discontinuità 
    \item il ripple nelle 2 bande è uguale 
    \item si vuole distribuire l'errore in frequenza 
    \item permettere di avere ripple diverso nelle bande 
\end{itemize} 

Quindi, considerando il filtro FIR progettato in equiripple (in figura quello celeste) ed il filtro ideale (in figura quello in rosso): 

\begin{figure}[h]
    \centering
    \includegraphics[scale = 0.35]{Filtro realizzato rispetto a quello ideale.png}
\end{figure} 

come il  caso IIR, dobbiamo considerare le specifiche come si vede nel seguente schema (con note): 

\begin{figure}[h]
    \centering
    \includegraphics[scale = 0.4]{filtro specifiche con note.png}
\end{figure} 

\newpage 

Questo tipo di filtro viene definito come filtro FIR equiripple perchè: 

\begin{itemize}
    \item l'errore viene distribuito nelle 2 bande uniformemente 
    \item a parità di M, si ottengono prestazioni migliori 
    \item il filtro è ottimo, nel senso che si minimizza l'errore massimo, cioè il valore massimo del ripple (questo significa ottimizzare)  
\end{itemize} 

Considerando sempre dei filtri a fase lineare generalizzata, come scritto in \ref{subsec: Sistemi causali con fase generalizzata lineare} \nameref{subsec: Sistemi causali con fase generalizzata lineare} a pag \pageref{subsec: Sistemi causali con fase generalizzata lineare}, 
perchè hanno fase lineare e ritardo di gruppo costante e la simmetria è garantita dalla fase lineare generalizzata. \newline 

Se consideriamo un filtro a fase lineare generalizzata di tipo I (ma potremmo anche considerare dei filtri di tipo II, III, IV, i passaggi sarebbero identici) 
non causale avremo che, per un M pari, il filtro sarà simmetrico ed a "fase zero", quindi: 

{
    \Large 
    \begin{equation}
        h \left[ n\right] 
        = 
        h \left[ -n\right]
        \text{  per M pari}
    \end{equation}
}

Allora, se poniamo: 

{
    \Large 
    \begin{equation}
        L = \dfrac{M}{2}
    \end{equation}
}

abbiamo che la trasformata di Fourier di $h \left[ n\right]$ vale: 

{
    \Large 
    \begin{equation}
        \setlength{\jot}{10pt} % affecting the line spacing in the environment
        \begin{split}
            H  \left(e^{\jmath \omega}\right) 
            &= 
            \sum_{- L}^{+ L} 
            h \left[ n\right]
            \cdot 
            e^{- \jmath \omega n} 
            \\
            &= 
            h \left[ 0\right] 
            + 
            \sum_{n = 1}^{+L} 
            2 \cdot h \left[ n\right] \cdot \cos \left(\omega n\right)
        \end{split}
    \end{equation}
}

L'espressione: 

{
    \Large 
    \begin{equation}
        H  \left(e^{\jmath \omega}\right) 
        = 
        h \left[ 0\right] 
        + 
        \sum_{n = 1}^{+L} 
        2 \cdot h \left[ n\right] \cdot \cos \left(\omega n\right)
    \end{equation}
} 

prende il nome di DCT (Discrete Cosine Transform). \newline 

\newpage 

\begin{tcolorbox}
    La DCT è un algoritmo impiegato principalmente per la compressione e viene utilizzato, in particolare, nella compressione JPEG. \newline 

    Se vuoi approfondire cosa è la DCT, ti consiglio queste pagine: 
    
    \begin{itemize}
        \item Trasformata discreta del coseno by Wikipedia \\\url{https://it.wikipedia.org/wiki/Trasformata_discreta_del_coseno} 
        \item Discrete Cosine Transform by Matlab\\\url{https://it.mathworks.com/help/images/discrete-cosine-transform.html}
    \end{itemize}    

    Invece, in formato video: 

    \begin{itemize}
        \item The Unreasonable Effectiveness of JPEG: A Signal Processing Approach by Reducible \\\url{https://youtu.be/0me3guauqOU?si=pnXh2Z05c2IDiPgi} 
        \item The Discrete Cosine Transform is magic! by theosib \\\url{https://youtu.be/QSBcLImmL4g?si=wC3RAIjKen9AaxB2} 
        \item JPEG DCT, Discrete Cosine Transform (JPEG Pt2)- Computerphile by Computerphile \\ \url{https://youtu.be/Q2aEzeMDHMA?si=5znc_gvA77ldVHwy}
    \end{itemize}
\end{tcolorbox}

Rispetto alla DFT, la DCT ha solo parti reali e non esiste il contributo delle parti immaginari. \newline 

Il problema della DCT è trovare una risposta impulsiva $h \left[ n\right]$ affinchè rispetti le specifiche: 

{
    \Large 
    \begin{equation}
        \begin{cases}
            \abs{1 - H \left(e^{\jmath \omega}\right)} \leq \delta_1 
            &\text{per } 
            0 \leq \abs{\omega} \leq \omega_p 
            \\
            \quad
            \\
            \abs{H \left(e^{\jmath \omega}\right)} \leq \delta_2 
            &\text{per } 
            \omega_s \leq \abs{\omega} \leq \pi 
        \end{cases}
    \end{equation}
}

Esistono vari metodi per trovare $h \left[ n\right]$: il più noto è l'algoritmo di Parks-McClellan (che verrà spiegato senza dimostrazione). \newline 

Fissati i seguenti valori: 

\begin{itemize}
    \item L 
    \item $\omega_p$ 
    \item $\omega_s$ 
    \item $\dfrac{\delta_1}{\delta_2}$
\end{itemize}

si ottiene: 

\begin{itemize}
    \item la funzione $h \left[ n\right]$ per i valori di $n = 0, \dots, L$ 
    \item i valori di $\delta_1$ e $\delta_2$
\end{itemize}

\newpage 

\subsection{Algoritmo di Parks-McClellan (1972)}

\footnote{ 
    Slide del prof | 10 - Progetto Filtri FIR 2 | pag 10 - 12\\
    Appunti | 2025-11-19 | pag 2 \\ 
    Latex | 2025-12-16 Mattino | pag 274  
}

Data la seguente funzione di $E \left(\omega\right)$, che prende il nome di funzione "errore di approssimazione": 

{
    \Large 
    \begin{equation}
        E \left(\omega\right) 
        = 
        W \left( \omega\right) 
        \cdot 
        \left[ 
            H_d \left(e^{\jmath \omega}\right) 
            - 
            H \left(e^{\jmath \omega}\right) 
        \right]
    \end{equation}
}

in cui: 

\begin{itemize}
    \item $W \left(\omega\right)$ è la funzione peso 
    \item $H_d \left(e^{\jmath \omega}\right)$ è la risposta in frequenza desiderata 
    \item $H \left(e^{\jmath \omega}\right)$ è la risposta del filtro in progetto 
\end{itemize}

quindi il valore di $E \left(\omega\right)$ non sarà costante, bensì varierà in base alla pulsazione $\omega$. \newline 

$W \left(\omega\right)$, $H_d \left(e^{\jmath \omega}\right)$ e $H \left(e^{\jmath \omega}\right)$ sono tutte funzioni reali e sono definite nella banda passante e nella banda oscura, 
ma non nella banda di transizione. \newline 

Si calcola $h \left[ n\right]$ per $n = 0, \dots, L$ affinchè: 

{
    \Large 
    \begin{equation}
        \max{\left| E \left(\omega\right) \right|_{\text{banda passante e banda oscura}}} 
    \end{equation}
}

sia minimo, utilizzando o il criterio MINMAX o il criterio di Chebyshev. \newline 

Nel caso di un filtro passa-basso precedente visto in \ref{subsubsec: Esempio di progettazione di filtri FIR} \nameref{subsubsec: Esempio di progettazione di filtri FIR} a pag \pageref{subsubsec: Esempio di progettazione di filtri FIR}, 
cioè un filtro di tipo I non causale con M pari ed a fase lineare generalizzata, quindi simmetrico ed a fase zero, in cui: 

{
    \Large 
    \begin{equation}
        K = \dfrac{\delta_1}{\delta_2}
    \end{equation}
} 

avremo che, grazie all'algoritmo di Parks-McClellan, la funzione "errore di approssimazione" $E \left(\omega\right)$ vale: 

{
    \Large 
    \begin{equation}
        E \left(\omega\right) 
        =
        \begin{cases}
        \setlength{\jot}{10pt} % affecting the line spacing in the environmen
        \dfrac{1}{K} 
        \cdot 
        \left[
            1 
            - 
            \left(
                h \left[ 0\right] 
                + 
                \sum_{n = 1}^{L} 
                2 \cdot h \left[ n\right] \cdot \cos \left(\omega n\right)
            \right)
        \right] 
        &\text{per } 0 \leq \abs{\omega} \leq \omega_p
        \\
        \quad
        \\
        1 
        \cdot 
        \left[
            0 
            - 
            \left(
                h \left[ 0\right] 
                + 
                \sum_{n = 1}^{L} 
                2 \cdot h \left[ n\right] \cdot \cos \left(\omega n\right)
            \right)
        \right] 
         &\text{per } \omega_s \leq \abs{\omega} \leq \pi
        \end{cases}
    \end{equation}
} 

Dalla formula di $E \left(\omega\right)$ possiamo svolgere le seguenti considerazioni: 

\begin{itemize}
    \item $E \left(\omega\right)$ dipende dal valore K, quindi sia da $\delta_1$ e $\delta_2$ perchè è presente il termine moltiplicativo $\dfrac{1}{K}$ nell'intervallo $0 \leq \abs{\omega} \leq \omega_p$, cioè nella banda passante   
    \item $E \left(\omega\right)$ non dipende dal valore K, quindi sia da $\delta_1$ e $\delta_2$ perchè non è presente il termine K nell'intervallo $\omega_s \leq \abs{\omega} \leq \pi$, cioè nella banda oscura 
    \item l'errore massimo pesato nelle due bande è per $\delta = \delta_2 $
    \item non c'è nessuna condizione per la banda di transizione, cioè per per $\omega_p \leq \abs{\omega}  \leq \omega_s$ 
\end{itemize} 

\newpage 

L'algoritmo di Parks-McClellan riformula il problema in termini di approssimazione polinomiale; 
in altre, parole, permette di calcolare la DFT come un polinomio, in particolare un polinomio di Chebyshev 
(che avevamo afforntato in precedenza con \ref{subsec: Filtri analogici di Chebyshev} \nameref{subsec: Filtri analogici di Chebyshev} a pag \pageref{subsec: Filtri analogici di Chebyshev}). \newline 

Tralasciando le dimostrazioni ed i passaggi matematici, 
grazie ai polinomio di Chebyshev, possiamo approssimare la funzione "errore di approssimazione" $E \left(\omega\right)$ come: 

{
    \Large 
    \begin{equation}
        E \left(\omega\right) 
        = 
        W \left(\omega\right) 
        \cdot 
        \left[ 
            H_d \left(e^{\jmath \omega}\right) 
            - 
            P \left(x\right)
        \right]_{x = \cos(\omega)}
    \end{equation}
}

dove $P \left(x\right)$ è il polinomio di Chebyshev. \newline 

Per massimizzare la funzione "errore di approssimazione", che a sua volta è approssimata, è necessario applicare il Teorema dell'alternanza. 

\newpage 

\subsubsection{Teorema dell'alternanza}

\footnote{ 
    Slide del prof | 10 - Progetto Filtri FIR 2 | pag 12 - 15\\
    Appunti | 2025-11-19 | pag 2  
}

\begin{tcolorbox}
    Questo teorema è da sapere per l'esame, senza il formalismo matematico. 
\end{tcolorbox}

Di seguito il teorema dell'alternanza. \newline 

Dato: 

{
    \Large 
    \begin{equation}
        E \left(\omega\right) 
        = 
        W \left(\omega\right) 
        \cdot 
        \left[ 
            H_d \left(e^{\jmath \omega}\right) 
            - 
            P \left(x\right)
        \right]_{x = \cos(\omega)}
    \end{equation}
} 

la condizione necessaria e sufficiente affinchè il $P \left(x\right)$ sia il solo polinomio che minimizza il valore massimo di $\abs{E \left(\omega\right)}$, 
è che $E \left(\omega\right)$ abbiamo almeno $\left(L + 2\right)$ alternanze nell'intervallo di definizione di $E \left(\omega\right)$, 
cioè esistano $\omega_1 < \omega_2 < \dots < \omega_{L+2}$ punti, tali che: 

{
    \Large 
    \begin{equation}
        E \left(\omega_i\right) 
        = 
        - E \left(\omega_{i+1}\right)
        = 
        \max \abs{E \left(\omega\right)} 
        \text{ per } 
        i = 1, \dots, \left(L+1\right)
    \end{equation}
}

Per il filtro passa-basso di tipo I, si ha dunque che ci devono essere almeno $\left(L + 2\right)$ punti di estremo negli intervalli chiusi $0 < \abs{\omega} < \omega_p$ e $\omega_s < \abs{\omega} < \pi$ 
come si vede nella seguente figura (con note): 

\begin{figure}[h]
    \centering
    \includegraphics[scale = 0.25]{filtro fir LPF Equiripple.png}
\end{figure} 

\newpage 

\subsubsection{Passi algoritmo di Parks-McClellan}

\footnote{ 
    Slide del prof | 10 - Progetto Filtri FIR 2 | pag 16\\
    Appunti | 2025-11-19 | pag 2  
}

\begin{tcolorbox}
    Questi passaggi bisogna saperli a memoria per l'esame. 
\end{tcolorbox}

Di seguito i passi dell'algoritmo di Parks-McClellan. \newline 

\begin{figure}[h]
    \centering
    \includegraphics[scale = 0.35]{passi algoritmo Parks-McClellan.png}
\end{figure} 

\newpage 

\subsection{Considerazioni finali dei filtri FIR Equiripple}

\footnote{ 
    Slide del prof | 10 - Progetto Filtri FIR 2 | pag 18\\
    Appunti | 2025-11-19 | pag 2  
}

Di seguito le considerazioni finali riguardo ai filtri FIR Equiripple: 

\begin{itemize}
    \item scegliere un valore di M adeguato in base all'utilizzo, se possibile un M più piccolo possibile 
    \item il caso "extraripple" porta a prestazioni migliori, ma si ottiene solo con particolari combinazioni di L, $\omega_p$, $\omega_s$ e K 
    \item per i filtri passa-basso e passa-alto vale la seguente relazione empirica di Kaiser: 
    {
        \Large 
        \begin{equation}
            M 
            \approx 
            \dfrac{
                10 \cdot \log_{10} \left(\delta_1 \cdot \delta_2\right) - 13 
            }{
                2.324 \cdot \left(\omega_s - \omega_p\right)
            }
        \end{equation}
    }
    \item per i filtri multi-banda, ci sono ancora almeno $\left(L + 2\right)$ estremi, ma non è necessario che siano solo nelle bande passanti e oscura
\end{itemize}

\newpage 
. 

\newpage 

